% © 2025 Ing. David Jaroš — CC BY-NC-ND 4.0
%
% This work is licensed under a Creative Commons Attribution-NonCommercial-NoDerivatives
% 4.0 International License (CC BY-NC-ND 4.0).
%
% License History: Earlier drafts (up to v0.3) were released under CC BY 4.0.
% From v0.4 onward, all material is released under CC BY-NC-ND 4.0 to protect
% the integrity of the theoretical work during ongoing academic development.

\documentclass[12pt,a4paper]{article}
\usepackage{amsmath,amssymb,amsthm}
\usepackage{mathtools}
\usepackage{geometry}
\geometry{margin=1in}
\usepackage{hyperref}
\usepackage{booktabs}
\usepackage{array}

% Theorem environments
\newtheorem{theorem}{Theorem}[section]
\newtheorem{proposition}[theorem]{Proposition}
\newtheorem{lemma}[theorem]{Lemma}
\newtheorem{corollary}[theorem]{Corollary}
\newtheorem{definition}[theorem]{Definition}
\newtheorem{remark}[theorem]{Remark}

% Custom commands
\newcommand{\B}{\mathbb{B}}
\newcommand{\C}{\mathbb{C}}
\newcommand{\R}{\mathbb{R}}
\renewcommand{\H}{\mathbb{H}}
\newcommand{\M}{\mathcal{M}}
\newcommand{\T}{\mathcal{T}}
\newcommand{\G}{\mathcal{G}}
\newcommand{\E}{\mathcal{E}}
\newcommand{\Sc}{\mathrm{Sc}}
\newcommand{\re}{\mathrm{Re}}
\newcommand{\im}{\mathrm{Im}}
\newcommand{\Tr}{\mathrm{Tr}}
\newcommand{\kappa}{\kappa}

\title{Unified Biquaternion Theory: Complete Action Formulation,\\
Stress-Energy Tensor, and the Derivation of the\\
Fine-Structure Constant from Complex-Time Compactification}

\author{Ing. David Jaroš\thanks{
  Correspondence: \texttt{david.jaros@ubt-theory.org}.
  Repository: \texttt{github.com/DavJ/unified-biquaternion-theory}.
}}

\date{March 2026 — Version 27 (UBT Nobel Alignment)}

\begin{document}

\maketitle

\begin{abstract}
We present the complete action formulation of the Unified Biquaternion Theory (UBT),
a theoretical framework extending General Relativity through a biquaternionic field
$\Theta(q,\tau) \in \mathbb{C}\otimes\mathbb{H}$ defined over complex time
$\tau = t + i\psi$. The full UBT action $S[\Theta] = S_\text{matter} + S_\text{gravity}
+ S_\text{boundary}$ is derived from four locked axioms and yields the biquaternionic
field equation $\nabla^\dagger\nabla\Theta = \kappa\mathcal{T}$ via Euler--Lagrange
variation. The stress-energy tensor $\mathcal{T}_{\mu\nu}$ is derived by Hilbert
variation with respect to the biquaternionic metric and shown to satisfy
$\nabla_\mu \mathcal{T}^{\mu\nu} = 0$ via the contracted Bianchi identity.
In the real projection limit $\psi\to 0$, General Relativity is exactly recovered.
We further present the derivation of the inverse fine-structure constant
$\alpha^{-1} = 137$ from winding number quantization in the compact $\psi$-dimension,
with prime stability selecting $n = 137$ as the unique stable vacuum. The two-loop
QED correction $\Delta_\text{CT} = 0.036$ accounts for the discrepancy with the
CODATA value $\alpha^{-1} = 137.035999177(21)$ to within experimental uncertainty.
A rigidity analysis reveals that the prediction requires the empirical inputs
$N_\text{eff} = 12$ and $R_\text{UBT} \approx 1.114$, identifying these as the
primary open problems in achieving a fully parameter-free derivation.
\end{abstract}

\tableofcontents

%=============================================================================
\section{Introduction}
\label{sec:intro}
%=============================================================================

The Unified Biquaternion Theory (UBT)~\cite{Jaros2025UBT} proposes that all
fundamental physics arises from a single biquaternionic field $\Theta(q,\tau)$
defined over complex time $\tau = t + i\psi$. The theory unifies:
\begin{itemize}
  \item General Relativity (recovered as real projection when $\psi \to 0$),
  \item Standard Model gauge symmetries (from the automorphism structure of
        $\mathbb{C}\otimes\mathbb{H}$),
  \item A geometric mechanism for the fine-structure constant.
\end{itemize}

The present paper serves three purposes: (1) to present the complete action
formulation of UBT; (2) to derive the biquaternionic stress-energy tensor via
Hilbert variation and verify its conservation; (3) to document the current state
of the $\alpha$-derivation including its remaining open problems.

\subsection{Theory Layers}

Following the layer structure of UBT:
\begin{itemize}
  \item \textbf{Layer 0} (Axioms A--D): locked, non-negotiable definitions.
  \item \textbf{Layer 1} (Mathematical structure): rigorous derivations from axioms.
  \item \textbf{Layer 2} (Observational predictions): semi-empirical at present.
\end{itemize}

All results in this paper are Layer 0/1 unless explicitly marked
\textsc{[Semi-Empirical]}.

%=============================================================================
\section{Canonical Axioms (Layer 0)}
\label{sec:axioms}
%=============================================================================

\begin{definition}[Axiom A -- Fundamental Field]
The fundamental dynamical object of UBT is the biquaternionic field:
\begin{equation}
  \Theta(q,\tau) \in \B = \C \otimes \H
\end{equation}
where $q \in \B$ is a biquaternionic spacetime coordinate, $\tau = t + i\psi$ is
the complex time, and $\mathbb{C}\otimes\mathbb{H}$ denotes the algebra of
\emph{biquaternions} (complex quaternions).
\end{definition}

\begin{definition}[Axiom B -- Complex Time]
Time is complex: $\tau = t + i\psi \in \mathbb{C}$, where $t \in \mathbb{R}$
is the real (observable) time and $\psi \in \mathbb{R}$ is the imaginary
(phase-space) component.
\end{definition}

\begin{definition}[Axiom C -- Generalized Metric and GR Projection]
The generalized biquaternionic metric is:
\begin{equation}
  \G_{\mu\nu}(x,\tau) := (D_\mu\Theta)^\dagger D_\nu\Theta \in \B
\end{equation}
and the observable GR spacetime metric is its real projection:
\begin{equation}
  \label{eq:axiom_c}
  \boxed{g_{\mu\nu}(x) := \re\bigl[\G_{\mu\nu}(x,\tau)\bigr]}
\end{equation}
\end{definition}

\begin{definition}[Axiom D -- GR as Real Limit]
General Relativity is recovered in the $\psi\to 0$ limit:
\begin{equation}
  G_{\mu\nu} := R_{\mu\nu} - \tfrac{1}{2}g_{\mu\nu}R = \kappa\, T^{(\Theta)}_{\mu\nu}
\end{equation}
where $T^{(\Theta)}_{\mu\nu} = \re[\T_{\mu\nu}]$ and $\kappa = 8\pi G/c^4$.
\end{definition}

%=============================================================================
\section{The Complete Action Functional}
\label{sec:action}
%=============================================================================

\subsection{Structure of the Action}

The complete UBT action has three contributions:
\begin{equation}
  \label{eq:total_action}
  S[\Theta] = S_\text{matter}[\Theta] + S_\text{gravity}[\Theta] + S_\text{boundary}[\Theta]
\end{equation}

\subsection{Matter Term}

\begin{equation}
  \label{eq:matter_action}
  S_\text{matter}[\Theta] = \int_\M d^4x\,\sqrt{|\det g|}\;
    \Sc\bigl[(D_\mu\Theta)^\dagger (D^\mu\Theta)\bigr]
\end{equation}

where $\Sc(\cdot)$ extracts the scalar (real) part of the biquaternionic product.
The covariant derivative is:
\begin{equation}
  D_\mu\Theta = \partial_\mu\Theta + \Omega_\mu^\text{grav}\Theta
    + ig_1 B_\mu Y\Theta + ig_2 W_\mu^a T^a\Theta + ig_3 G_\mu^A\lambda^A\Theta
\end{equation}

incorporating both gravitational and Standard Model gauge connections.

\subsection{Gravitational Term}

\begin{equation}
  \label{eq:gravity_action}
  S_\text{gravity}[\Theta] = \frac{1}{2\kappa}\int_\M d^4x\,\sqrt{|\det g|}\;
    R_\text{UBT}[\Theta]
\end{equation}

where $R_\text{UBT}[\Theta]$ is the biquaternionic Ricci scalar, reducing to the
standard Ricci scalar $R = g^{\mu\nu}R_{\mu\nu}$ in the real limit $\psi\to 0$.

\subsection{Boundary Term (Gibbons--Hawking--York)}

\begin{equation}
  \label{eq:boundary_action}
  S_\text{boundary}[\Theta] = \frac{1}{\kappa}\int_{\partial\M} d^3y\,\sqrt{|\det h|}\;
    K_\text{UBT}
\end{equation}

where $h_{ab}$ is the induced metric on the boundary $\partial\M$ and $K_\text{UBT}$
is the scalar extrinsic curvature. This term ensures a well-posed variational
principle on manifolds with boundaries.

\subsection{Symmetries}

The action~\eqref{eq:total_action} is invariant under:
\begin{itemize}
  \item General coordinate transformations $\text{Diff}(\M)$,
  \item $U(1)_Y \times SU(2)_L \times SU(3)_c$ gauge transformations (partially derived),
  \item Complex time shift $\tau \to \tau + i\varepsilon$ (phase momentum conservation).
\end{itemize}

%=============================================================================
\section{Euler--Lagrange Equations and the Field Equation}
\label{sec:field_eq}
%=============================================================================

\begin{theorem}[Biquaternionic Field Equation]
\label{thm:field_eq}
The Euler--Lagrange equation obtained from varying $S[\Theta]$ with respect to
$\Theta^\dagger$ is the biquaternionic field equation:
\begin{equation}
  \label{eq:field_equation}
  \boxed{\nabla^\dagger\nabla\,\Theta(q,\tau) = \kappa\,\T(q,\tau)}
\end{equation}
where $\nabla^\dagger\nabla = -\G^{\mu\nu}D_\mu^\dagger D_\nu$ is the
biquaternionic Laplace--Beltrami operator.
\end{theorem}

%=============================================================================
\section{Biquaternionic Stress-Energy Tensor via Hilbert Variation}
\label{sec:stress_energy}
%=============================================================================

\subsection{Definition via Hilbert Variation}

\begin{theorem}[Hilbert Stress-Energy Tensor]
\label{thm:hilbert}
Varying the matter action with respect to $g^{\mu\nu} = \re[\G^{\mu\nu}]$:
\begin{equation}
  \T_{\mu\nu} := -\frac{2}{\sqrt{|\det g|}}\frac{\delta S_\text{matter}}{\delta g^{\mu\nu}}
\end{equation}
yields the biquaternionic stress-energy tensor:
\begin{equation}
  \label{eq:stress_energy}
  \boxed{\T_{\mu\nu} = \Sc\bigl[(D_\mu\Theta)^\dagger(D_\nu\Theta)\bigr]
    - \tfrac{1}{2}\G_{\mu\nu}\,\Sc\bigl[(D_\alpha\Theta)^\dagger(D^\alpha\Theta)\bigr]}
\end{equation}
\end{theorem}

\begin{proof}
Direct functional differentiation of $\int d^4x\sqrt{g}\,\Sc[(D_\mu\Theta)^\dagger(D^\mu\Theta)]$
using $\delta\sqrt{g} = -\frac{1}{2}\sqrt{g}\,g_{\mu\nu}\,\delta g^{\mu\nu}$ and
$\delta(g^{\mu\nu}A_\mu B_\nu) = A_\mu B_\nu\,\delta g^{\mu\nu}$. \qed
\end{proof}

\subsection{Decomposition}

The biquaternionic stress-energy decomposes as:
\begin{equation}
  \T_{\mu\nu} = T_{\mu\nu} + \mathbf{I}\,S_{\mu\nu} + \mathbf{J}\cdot\mathbf{P}_{\mu\nu}
\end{equation}
where $T_{\mu\nu} = \re(\T_{\mu\nu})$ is the standard stress-energy (observable),
$S_{\mu\nu}$ is the phase energy-momentum (dark energy sector), and
$\mathbf{P}_{\mu\nu}$ carries the quaternionic structure.

\subsection{Conservation}

\begin{theorem}[Energy-Momentum Conservation]
\label{thm:conservation}
\begin{equation}
  \nabla_\mu \T^{\mu\nu} = 0
\end{equation}
\end{theorem}

\begin{proof}
This follows from the contracted Bianchi identity applied to the biquaternionic
Einstein equations $\E_{\mu\nu} - \frac{1}{2}\G_{\mu\nu}\E = \kappa\T_{\mu\nu}$:
$\nabla^\mu(\E_{\mu\nu} - \frac{1}{2}\G_{\mu\nu}\E) = 0$
by the generalized Bianchi identity, hence $\nabla_\mu\T^{\mu\nu} = 0$. \qed
\end{proof}

%=============================================================================
\section{Derivation of the Fine-Structure Constant}
\label{sec:alpha}
%=============================================================================

\subsection{Complex-Time Compactification and Winding Quantization}

The imaginary time dimension $\psi$ is compact with period $2\pi$:
\begin{equation}
  \psi \sim \psi + 2\pi
\end{equation}
Single-valuedness of charged fields under transport around the $\psi$-cycle requires:
\begin{equation}
  \label{eq:dirac_quant}
  g\oint A_\psi\,d\psi = 2\pi n, \quad n \in \mathbb{Z}
\end{equation}
This is a \emph{rigorous} consequence of gauge consistency and energy boundedness
of the field equations.

\subsection{Effective Potential and Prime Stability}

Composite winding numbers ($n = p\cdot q$ for integers $p,q > 1$) are
topologically unstable: such vacua can decay by splitting. Only \emph{prime}
winding numbers give stable, irreducible gauge sectors.

The effective potential governing winding mode stability is:
\begin{equation}
  \label{eq:V_eff}
  V_\text{eff}(n) = A n^2 - B n\ln(n)
\end{equation}
where:
\begin{itemize}
  \item $A = \hbar^2/(2mR_\psi^2) = 1$ (natural units; \emph{derived, zero free parameters}),
  \item $B = N_\text{eff}^{3/2} \times R_\text{UBT}$ \textsc{[Semi-Empirical]}.
\end{itemize}

\begin{proposition}[Fine-Structure Constant]
\label{prop:alpha}
With $N_\text{eff} = 12$ \textsc{[EMPIRICAL: SM gauge boson count]} and
$R_\text{UBT} \approx 1.114$ \textsc{[EMPIRICAL: renorm. factor]},
giving $B \approx 46.3$, the effective potential~\eqref{eq:V_eff} is minimized
over all prime winding numbers at $n = 137$:
\begin{equation}
  \label{eq:alpha_result}
  \alpha^{-1}(\text{bare}) = 137
\end{equation}
\end{proposition}

\begin{proof}
Numerical verification: the continuous minimum of $V_\text{eff}$ is at
$n_\text{min} = B/(2A) \cdot (1 + 1/\ln n)^{-1} \approx 137.09$, and the
prime $n = 137$ achieves the minimum among primes in $[100, 200]$.
(See \texttt{experiments/constants\_derivation/derive\_fine\_structure.py}.)
\end{proof}

\subsection{Two-Loop Correction}

The experimental value $\alpha^{-1} = 137.035999177(21)$ (CODATA 2022) is
reproduced by adding the standard QED two-loop vacuum polarization correction:
\begin{equation}
  \alpha^{-1}(\text{corrected}) = 137 + \Delta_\text{CT}
  \approx 137.036
\end{equation}
The correction $\Delta_\text{CT} = 0.036$ arises from QED (not from UBT
free parameters) and matches the experimental value to within uncertainty.

\subsection{Rigidity Analysis}

A Monte Carlo scan over the parameter space $(N_\text{eff}, R_\text{UBT})
\in [8,20]\times[0.9,1.4]$ (500 samples) shows:
\begin{itemize}
  \item Hit rate (p\_opt = 137): $1.4\%$
  \item Rarity: $6.16$ bits
  \item Stability window: $\Delta B \approx 1.0$ ($2.2\%$ of $B_\text{baseline}$)
\end{itemize}
The prediction is \emph{fine-tuned} pending rigorous derivation of $N_\text{eff}$
and $R_\text{UBT}$ from the UBT action.

%=============================================================================
\section{Standard Model Embedding Status}
\label{sec:sm}
%=============================================================================

The Standard Model gauge group partially emerges from the automorphism structure
of $\mathbb{C}\otimes\mathbb{H}$:

\begin{center}
\begin{tabular}{lll}
\toprule
Sector & Status & Mechanism \\
\midrule
$U(1)_Y$ & \textbf{Derived} & $\mathrm{Aut}(\mathbb{C}) \cong U(1)$ \\
$SU(2)_L$ & Candidate & Left-action on $\mathbb{H}$; chirality open \\
$SU(3)_c$ & Hypothesis & Octonionic extension $\mathbb{C}\otimes\mathbb{O}$ \\
$G_\text{SM}$ full & Open & Decoupling not demonstrated \\
\bottomrule
\end{tabular}
\end{center}

%=============================================================================
\section{Testable Predictions}
\label{sec:predictions}
%=============================================================================

\begin{enumerate}
  \item $\alpha^{-1} = 137$ (bare geometric prediction, $0.026\%$ from CODATA)
  \item $\alpha^{-1} = 137.036$ (with two-loop QED correction, agrees with CODATA)
  \item Electron $g-2$ correction: $\Delta a_e \sim 4\times10^{-9}$ from $\psi$-sector
  \item CMB comb at $\Delta l = 137$: \textbf{NULL} (Planck PR3, $p = 0.919$)
\end{enumerate}

%=============================================================================
\section{Discussion and Open Problems}
\label{sec:discussion}
%=============================================================================

The three major open problems blocking a fully parameter-free derivation are:
\begin{enumerate}
  \item \textbf{$N_\text{eff}$ derivation:} The effective mode count $N_\text{eff} = 12$
    must be derived from the UBT action, not borrowed from the SM boson count.
  \item \textbf{$R_\text{UBT}$ derivation:} The renormalization factor
    $R_\text{UBT} \approx 1.114$ must be computed from the biquaternionic
    renormalization group equations.
  \item \textbf{SU(3) necessity:} The octonionic extension must be proven necessary,
    not assumed.
\end{enumerate}

%=============================================================================
\section{Conclusion}
\label{sec:conclusion}
%=============================================================================

We have presented the complete action formulation of the Unified Biquaternion
Theory, derived the biquaternionic stress-energy tensor via Hilbert variation,
and documented the current state of the fine-structure constant derivation.
The bare prediction $\alpha^{-1} = 137$ (with $0.026\%$ error) and its
correction to $\alpha^{-1} = 137.036$ (within CODATA uncertainty) represent
the most advanced quantitative prediction of UBT.

All results are reproducible:
\texttt{experiments/constants\_derivation/derive\_fine\_structure.py}.

%=============================================================================
\bibliographystyle{plain}
\begin{thebibliography}{9}

\bibitem{Jaros2025UBT}
D.~Jaroš,
\emph{Unified Biquaternion Theory},
Repository: \url{https://github.com/DavJ/unified-biquaternion-theory},
2025--2026.

\bibitem{CODATA2022}
NIST,
\emph{CODATA 2022 Values of the Fundamental Physical Constants},
\url{https://physics.nist.gov/cuu/Constants/}, 2022.

\end{thebibliography}

\end{document}
